% ============================================================
% Chapter I「导引」§I.1–§I.4（PDF p.7–p.22）
% 底本：Fargues–Scholze, Geometrization of the local Langlands
%       correspondence, arXiv:2102.13459v4，印刷页 = PDF 页，偏移 0。
% 体例：逐段对照；每页起点 \srcpage{N}；译注 \yizhu{注文}{来源}。
% ============================================================

\chapter{导引}% 原文 Chapter I: Introduction

\section{局部 Langlands 对应}% 原文 I.1. The local Langlands correspondence

\leavevmode\srcpageinl{7}局部 Langlands 对应旨在描述 $G(E)$ 的不可约光滑表示 $\pi$，其中 $G$ 是局部域 $E$ 上的约化群。在下文另行说明之前，我们假定 $G$ 是分裂的，以简化处境；正文处理一般约化群 $G$ 的情形，只需作微小的改动。

当 $E$ 是阿基米德域，即 $E=\RR$ 或 $E=\CC$ 时，这是 Langlands 经典工作 [Lan89] 的主题。基于 Harish-Chandra 的工作（例如可参看 [HC66]），Langlands 给每个 $\pi$ 配一条 \emph{$L$-参数}，即一个连续同态
\[
  \varphi_\pi : W_E \longrightarrow \widehat{G}(\CC)
\]
其中 $W_E$ 是 $E=\RR,\CC$ 的 Weil 群（$W_\CC=\CC^\times$；相应地 $W_\RR$ 是非分裂扩张 $1\to W_\CC\to W_\RR\to\Gal(\CC/\RR)\to 1$），而 $\widehat{G}$ 是朗兰兹对偶群%
\yizhu{Langlands dual group 本书统一译作「朗兰兹对偶群」（术语表另收别名「Langlands 对偶群」）：它是 $\ZZ$ 上分裂的约化群，其根数据与 $G$ 的根数据对偶。}{I：本工程 glossary.tsv；A：原文 §I.1 第 2 段}。
它是 $\ZZ$ 上分裂的约化群，其根数据对偶于 $G$ 的根数据。映射 $\pi\mapsto\varphi_\pi$ 的纤维有限，且有大量工作致力于把这些纤维，即所谓的 \emph{$L$-包}，显式地刻画出来\yizhu{「$L$-包」（$L$-packet）指同一 $L$-参数 $\varphi_\pi$ 之下的不可约表示所成的纤维；它的大小正是局部 Langlands 对应偏离单射的程度。}{A：原文 §I.1 第 2 段与 §I.9}
。若 $G=\mathrm{GL}_n$，则映射 $\pi\mapsto\varphi_\pi$ 本质上是双射%
\yizhu{原文的「本质上是双射」须与上一句合看：$\mathrm{GL}_n$ 的 $L$-包是单点，而参数一侧只计半单 $L$-参数。原文 p.36（§I.9）明确写道：对 $\mathrm{GL}_n$ 与一般的 $\pi$，$\varphi_\pi$ 即通常所谓的半单 $L$-参数。}{A：原文 p.36}。

贯穿本文，我们假定 $E$ 是非阿基米德的，剩余特征 $p>0$，剩余域为 $\FF_q$。Langlands 猜想，仍可自然地给 $G(E)$ 的任一不可约光滑表示 $\pi$ 配一条 $L$-参数
\[
  \varphi_\pi : W_E \longrightarrow \widehat{G}(\CC)
\]
在非阿基米德情形，$W_E$ 是绝对 Galois 群 $\Gal(\ol{E}|E)$ 的稠密子群，由 $\ZZ\subset\Gal(\ol{\FF}_q|\FF_q)$ 的原像给出，其中后者由 Frobenius $x\mapsto x^q$ 生成。这就引出问题：这样一个参数应当从何而来。特别地，

(1) Weil 群 $W_E$ 与 $G(E)$ 的表示论如何相关？

(2) 朗兰兹对偶群 $\widehat{G}$ 从何而来？

本文的目标是给出参数 $\varphi_\pi$ 的一个自然构造（它只依赖于同构 $\CC\cong\ol{\QQ}_\ell$ 的选取），并在此过程中回答上面两个问题 (1) 与 (2)。

\section{总体图景}% 原文 I.2. The big picture

\leavevmode\srcpage{8}在代数几何中，任一环 $A$ 对应一个空间 $\Spec A$。我们考察的出发点是仔细审视与 $E$ 相伴的空间 $\Spec E$。\footnote{不消说，下面的叙述与这些思想的历史发展毫无关系；一如既往，实际的历史路径要迂回得多。我们将在下文 §I.11 讨论我们的一部分原始动机。} 注意，群 $G(E)$ 是 $\Spec E$ 上平凡 $G$-主丛的自同构群，而 $E$ 的 Weil 群本质上是 $E$ 的绝对 Galois 群，即 $\Spec E$ 的（étale）基本群。于是 $G(E)$ 与 $\Spec E$ 上的凝聚信息（特别是 $G$-主丛）相关，而 $W_E$ 与 $\Spec E$ 上的 étale 信息相关。此外，$G$-主丛的视角是个好视角：对一般的群 $G$，$\Spec E$ 上可以存在非平凡的 $G$-主丛 $\mathcal{E}$，其自同构群就是 Vogan [Vog93] 所谓的纯内形式\yizhu{「纯内形式」（pure inner form）指由一个 $G$-主丛 $\mathcal{E}$ 给出的内形式 $\Aut(\mathcal{E})$；Vogan 引入它是为了把 $L$-包的纤维一并参数化。}{A：原文 §I.2；Vogan, \emph{The local Langlands conjecture}（[Vog93]）}
。Vogan 意识到，从局部 Langlands 对应、特别是从 $\pi\mapsto\varphi_\pi$ 的纤维的参数化的角度看，把所有纯内形式放在一起考虑是有益的；特别地，他由此能对（拟）分裂群的纯内形式表述出局部 Langlands 猜想的一个精确形式（把 $\pi\mapsto\varphi_\pi$ 的纤维也考虑在内）。所有纯内形式合起来，来自考察 $\Spec E$ 上所有 $G$-丛构成的 groupoid：它由
\[
  [\ast/G](\Spec E) \;=\; \bigsqcup_{[\alpha]\in H^1_{\text{ét}}(\Spec E,G)} [\ast/G_\alpha(E)]
\]
给出，其中 $H^1_{\text{ét}}(\Spec E,G)$ 是 $\Spec E$ 上 $G$-主丛的同构类之集，而 $G_\alpha$ 是 $G$ 的相应纯内形式。我们还要指出，$G(E)$ 的表示等价于 $[\ast/G(E)]$ 上的层（若 $G(E)$ 是离散群，这是同义反复；在目前光滑表示的语境中，取正确的「层」概念，它同样成立），因此
\[
  [\ast/G](\Spec E) \;=\; \bigsqcup_{[\alpha]\in H^1_{\text{ét}}(\Spec E,G)} [\ast/G_\alpha(E)]
\]
上的层等价于表示组 $(\pi_\alpha)_{[\alpha]}$，其中 $\pi_\alpha$ 是 $G_\alpha(E)$ 的表示。\footnote{把 $[\ast/G(E)]$ 换成 $[\ast/G](\Spec E)$ 的这一点，Bernstein 也曾强调过。}

考察对应的 étale 一侧，我们注意到，局部 Langlands 对应所显现的是 $E$ 的 Weil 群 $W_E$，而不是它的绝对 Galois 群 $\Gal(\ol{E}|E)$。回忆 $W_E\subset\Gal(\ol{E}|E)$ 是由包含 $\ZZ\subset\Gal(\ol{\FF}_q|\FF_q)\cong\hat{\ZZ}$ 的原像给出的稠密子群，其中 $\Gal(\ol{\FF}_q|\FF_q)$ 由它的 Frobenius 态射 $x\mapsto x^q$ 生成。在几何的层面上，这一更换对应于把 $\FF_q$ 上的概形 $X$ 换成（形式）商 $X_{\ol{\FF}_q}/\mathrm{Frob}$。

在函数域情形 $E=\FF_q((t))$，我们因此被引向把 $\Spec E$ 换成 $\Spec\breve{E}/\varphi^\ZZ$，其中 $\breve{E}=\ol{\FF}_q((t))$。其实一般情形也可照做，取 $\breve{E}$ 为 $E$ 的极大无分歧扩域的完备化。对 $\pi_1$ 作自然的定义后，就有 $\pi_1(\Spec(\breve{E})/\varphi^\ZZ)=W_E$——例如 $\Spec\ol{E}\to\Spec(\breve{E})/\varphi^\ZZ$ 是一个 $W_E$-主丛，其中 $\ol{E}$ 是可分闭包。

让我们分析这一替换在对应的另一侧意味着什么：考察 $\Spec\breve{E}/\varphi^\ZZ$ 的凝聚理论，人们被引向研究带有 $\varphi$-线性自同构 $\sigma$ 的 $\breve{E}$-向量空间 $V$。这就是所谓的等晶体范畴 $\Isoc_E$\yizhu{isocrystal 本书统一译作「等晶体」（字面对应「同晶」义，取「同构类意义下的晶体」之意）；它是 $\breve{E}$ 上带 $\varphi$-线性自同构的向量空间所成范畴。}{I：译名取舍为译者说明；原文 §I.2 只给出该词所指的对象，未论译名}
。等晶体范畴比 $E$-向量空间范畴丰富得多，而后者全忠实地嵌入其中。具体地说，由 \srcpageinl{9}Dieudonné–Manin 分类，等晶体范畴 $\Isoc_E$ 是半单的，对每个有理数 $\lambda\in\QQ$ 恰有一个单对象 $V_\lambda$。$V_\lambda$ 的自同态代数由 Brauer 不变量为 $\lambda\in\QQ/\ZZ$ 的中心单 $E$-代数 $D_\lambda$ 给出。于是有范畴等价
\[
  \Isoc_E \;\cong\; \bigoplus_{\lambda\in\QQ} \Vect_{D_\lambda}\otimes V_\lambda .
\]
这里，若把 $\lambda$ 写成 $\lambda=\frac{s}{r}$（$s,r$ 互素，$r>0$），则 $V_\lambda\cong\breve{E}^r$ 秩为 $r$，而 $\sigma$ 由矩阵
\[
  \begin{pmatrix}
    0 & 1 & 0 & \cdots & 0\\
      & 0 & 1 & & 0\\
      & & \ddots & \ddots & \vdots\\
    0 & & & 0 & 1\\
    \pi^s & 0 & 0 & \cdots & 0
  \end{pmatrix}
\]
给出，其中 $\pi\in E$ 是一致化子。

上面我们考虑的是 $\Spec E$ 上的 $G$-主丛，故现在应考察 $\Isoc_E$ 中的 $G$-主丛。它们称为 $G$-等晶体，已由 Kottwitz [Kot85]、[Kot97] 广泛研究。其研究最初的动机是等晶体与 $p$-可除群的关系，以及相应地 $G$-等晶体与 Shimura 簇特殊纤维的关系（后者参数化带 $G$-结构的阿贝尔簇，从而参数化带 $G$-结构的 $p$-可除群）。传统上，$G$-等晶体之集记作 $B(E,G)$，而对 $b\in B(E,G)$ 我们记 $\mathcal{E}_b$ 为相应的 $G$-等晶体。特别地，Kottwitz 分离出了基本的 $G$-等晶体这一类；对 $G=\mathrm{GL}_n$，$G$-等晶体就是秩 $n$ 的等晶体，它是基本的当且仅当它只有一条斜率 $\lambda$。存在单射 $H^1_{\text{ét}}(\Spec E,G)\hookrightarrow B(E,G)$，因为 $\Spec E$ 上任一 $G$-主丛都「拉回」为 $\Isoc_E$ 中的一个 $G$-主丛\yizhu{单射性来自 $E$-向量空间范畴到 $\Isoc_E$ 的全忠实嵌入（见上一页末）：若两个 $G$-主丛拉回后同构，则它们本身已同构。}{I：此步为译者补证，原文只说「由全忠实性立得」；A：原文 p.8 末段与 p.9}
；像落在 $B(E,G)_{\mathrm{basic}}$ 中。对任意 $b\in B(E,G)_{\mathrm{basic}}$，$\mathcal{E}_b$ 的自同构群是 $G$ 的一个内形式 $G_b$；$G$ 的这类内形式之集称为 $G$ 的扩充纯内形式。注意对 $G=\mathrm{GL}_n$，$G$ 没有非平凡的纯内形式，但 $G$ 的所有内形式都是扩充纯内形式，这恰是因为所有中心单 $E$-代数都作为某个斜率 $\lambda$ 的 $H_\lambda$ 出现。更一般地，若 $G$ 的中心连通，则 $G$ 的所有内形式都可实现为扩充纯内形式。Kaletha [Kal14] 把 Vogan 关于纯内形式的结果推广到扩充纯内形式，对所有扩充纯内形式给出了局部 Langlands 对应的精确形式（描述 $\pi\mapsto\varphi_\pi$ 的纤维），从而表明 $G$-等晶体从纯粹表示论的观点看也是有益的。我们实际上将在下文论证：最好把所有 $b\in B(E,G)$ 的 $G_b$ 都纳入，而不只是基本的 $b$；此时得到的自同构群 $G_b$ 是 $G$ 的 Levi 子群的内形式。于是我们被引向考虑 $\Isoc_E$ 中 $G$-主丛的 groupoid
\[
  G\text{-}\Isoc \;\cong\; \bigsqcup_{[b]\in B(E,G)} [\ast/G_b(E)] .
\]
其上的层于是就是表示组 $(\pi_b)_{[b]\in B(E,G)}$，其中 $\pi_b$ 是 $G_b(E)$ 的表示。局部 Langlands 猜想，包括它在过渡到内形式与 Levi 子群方面所预期的函子性行为，此时仍然预测：对任一不可约层 $\mathcal{F}$——它必然由某个 $b\in B(E,G)$ 的不可约表示 $\pi_b$ 给出——都可以配一条 $L$-参数
\[
  \varphi_{\mathcal{F}} : W_E \longrightarrow \widehat{G}(\CC).
\]

\leavevmode\srcpageinl{10}要进一步，我们需要把几何引入图景：事实上，正是通过几何，$\Spec\breve{E}/\varphi^\ZZ$ 上的 $G$-主丛（其 groupoid 上的层）才会与 $\Spec\breve{E}/\varphi^\ZZ$ 的基本群 $W_E$ 联系起来。关键的想法是研究 $\Spec\breve{E}/\varphi^\ZZ$ 上 $G$-主丛的模叠。

要定义这样一个模叠，有几种可尝试的方式。最朴素的或许是下面这种。范畴 $\Isoc_E$ 是 $E$-线性的。于是对任意 $E$-代数 $A$，我们可以考虑 $\Isoc_E\otimes_E A$ 中的 $G$-主丛。所得的模叠实际上将是
\[
  \bigsqcup_{b\in B(E,G)} [\ast/G_b],
\]
即 $E$ 上的 Artin 叠，由各代数群 $G_b$ 的分类叠之不交并给出。这一视角对把诸 $G_b$ 定义为代数群实际上是很有用的。然而，它对进一步几何化处境这一目标并无帮助。这是因为 $[\ast/G_b]$ 上的层是代数群 $G_b$ 的表示，而我们所关心的是局部 profinite 群 $G_b(E)$ 的表示。

更好的视角是把 $\FF_q$ 的选取视为辅助性的，并用一般的 $\FF_q$-代数 $R$ 来代替它。在等特征情形，我们于是可把 $\breve{E}=\FF_q((t))$ 换成 $R((t))$。这带有作用在 $R$ 上的 Frobenius $\varphi=\varphi_R$。为过渡到商 $\Spec R((t))/\varphi^\ZZ$，需要假定 $R$ 的 Frobenius 是自同构，即 $R$ 是完满的。（限制到完满的 $R$ 在混合特征情形将变得更为关键。就考虑 $\ell$-进层而言，过渡到完满环是无害的，因为概形 $X$ 上与其完满化上的 étale 层自然地等价\yizhu{完满化的射影 $X_{\mathrm{perf}}\to X$ 是万有同胚，故两者的 étale 站点给出相同的层范畴；这正是原文所说「自然地等价」的含义。}{I：此解释为译者补证；A：原文 §I.2 仅有「自然地等价」一句}
。）我们于是被引向完满 $\FF_q$-代数上的模叠
\[
  G\text{-}\mathcal{I}\mathrm{soc} : \{\text{perfect } \FF_q\text{-algebras}\} \to \{\text{groupoids}\} : R \mapsto \{G\text{-torsors on } \Spec R((t))/\varphi^{\ZZ}\}.
\]
它们也称为完满概形 $\Spec R$ 上 $G$-等晶体的族。（注意 $G$-$\mathcal{I}$soc 中花体的那个 $I$，以区别于 groupoid $G$-$\Isoc$。）

这一构造可以推广到混合特征情形。事实上，若 $R$ 是完满 $\FF_q$-代数，则 $R[[t]]$ 的类比是唯一的 $\pi$-进完备平坦 $O_E$-代数 $\widetilde{R}$，使得 $\widetilde{R}/\pi=R$；用 $p$-典型 Witt 向量 $W(R)$ 或分歧 Witt 向量 $W_{O_E}(R)$ 显式写出，就是 $\widetilde{R}=W_{O_E}(R)=W(R)\otimes_{W(\FF_q)}O_E$。因此，若 $E$ 是混合特征的，我们定义
\[
  G\text{-}\mathcal{I}\mathrm{soc} : \{\text{perfect } \FF_q\text{-algebras}\} \to \{\text{groupoids}\} : R \mapsto \{G\text{-torsors on } \Spec(W_{O_E}(R)[\tfrac{1}{\pi}])/\varphi^{\ZZ}\}.
\]

本文不会用到叠 $G$-$\mathcal{I}$soc。不过，近来它已被若干人强调，其中隐含地由 Genestier–V. Lafforgue [GL17] 提出，明确地由 Gaitsgory [Gai16, §4.2] 与 Zhu [Zhu20] 提出；人们可以期望本文的结果在 $G$-$\mathcal{I}$soc 的语言下也有平行的表述，故让我们在本文导引中进一步分析它。它常以下面稍有不同的形式被定义。即：由 Anschütz 近来的结果 [Ans22]，在 $R$ 上 v-局部地，$R((t))$ 上（resp.\ $W_{O_E}(R)[\tfrac1\pi]$ 上）的任何 $G$-主丛都是平凡的。选定这样一个平凡化之后，一个 $G$-等晶体的族就由 $LG(R)$ 的某个元素给出，其中我们如下定义环路群
\[
  LG(R)=G(R((t)))\qquad(\text{resp. } LG(R)=G(W_{O_E}(R)[\tfrac{1}{\pi}])).
\]
改变 $G$-主丛的平凡化相当于在 $LG$ 上作 $\sigma$-共轭，故作为 v-叠
\[
  G\text{-}\mathcal{I}\mathrm{soc} = LG/_{\mathrm{Ad},\sigma}LG
\]
是 $LG$ 在 $\sigma$-共轭作用下按 $LG$ 取的商。

\leavevmode\srcpage{11}叠 $G$-$\mathcal{I}$soc 是可以分析的。更确切地说，我们有下面的结果。\footnote{这个结果似乎为专家所熟知，但我们并不知道一份完整的参考文献。关于 v-下降（乃至 arc-下降），见 [Iva23, Lemma 5.9]。该分层本质上是在 [RR96] 中构造的；Kottwitz 映射的局部常性在本文推论 III.2.8 中一般地证明了。在某些情形下，各层的识别见 [CS17, Proposition 4.3.13]，一般情形见 [HK22, Theorem 1.4]。}

\begin{theorem}% 原编号 I.2.1\label{thm:I.2.1}
prestack $G$-$\mathcal{I}$soc 对完满 $\FF_q$-代数上的 v-拓扑是叠。它容许一个由局部闭子叠构成的分层
\[
  G\text{-}\mathcal{I}\mathrm{soc}^b \subset G\text{-}\mathcal{I}\mathrm{soc}
\]
（$b\in B(E,G)$），由在每个几何点处 $G$-等晶体同构于 $\mathcal{E}_b$ 的轨迹组成。此外，每个层
\[
  G\text{-}\mathcal{I}\mathrm{soc}^b \cong [\ast/G_b(E)]
\]
都是局部 profinite 群 $G_b(E)$ 的分类叠。
\end{theorem}

环路群 $LG$ 是 ind-(无穷维完满概形)，故上述表述
\[
  G\text{-}\mathcal{I}\mathrm{soc} = LG/_{\mathrm{Ad},\sigma}LG
\]
具有极度无限的性质。我们猜测，问题并不出在这个表述上，而是这个叠本身没有好的有限性；特别地请注意，各层看上去维数都是 $0$，却容许非平凡的特殊化关系。尽管如此，Xiao–Zhu（见 [XZ17]、[Zhu20]）仍能定义 $G$-$\mathcal{I}$soc 上 $\ell$-进层的范畴 $D(G\text{-}\mathcal{I}\mathrm{soc},\ol{\QQ}_\ell)$，它容许半正交分解，其各项为各个 $D(G\text{-}\mathcal{I}\mathrm{soc}^b,\ol{\QQ}_\ell)$。每个 $D(G\text{-}\mathcal{I}\mathrm{soc}^b,\ol{\QQ}_\ell)\cong D([\ast/G_b(E)],\ol{\QQ}_\ell)$ 都等价于 $G_b(E)$ 的光滑表示（在 $\ol{\QQ}_\ell$-向量空间上）所成范畴的导出范畴。这里一如既往，我们必须固定一个辅助素数 $\ell\neq p$ 以及同构 $\CC\cong\ol{\QQ}_\ell$。

至此我们已经定义了一个叠 $G$-$\mathcal{I}$soc，并带有闭浸
\[
  i : [\ast/G(E)] \cong G\text{-}\mathcal{I}\mathrm{soc}^1 \subset G\text{-}\mathcal{I}\mathrm{soc},
\]
从而实现了全忠实嵌入
\[
  i_\ast : D(G(E),\ol{\QQ}_\ell) \hookrightarrow D(G\text{-}\mathcal{I}\mathrm{soc},\ol{\QQ}_\ell),
\]
即把 $G(E)$ 的光滑表示的导出范畴嵌入到 $G$-$\mathcal{I}$soc 上 $\ol{\QQ}_\ell$-层之导出范畴中。正是以这种方式，我们「几何化了 $G(E)$ 的表示论」。

我们所需的关键额外结构是 Hecke 算子\yizhu{「Hecke 算子」是几何 Langlands 纲领的核心结构：它由沿曲线作修正（modification）给出，并同时让 Weil 群与朗兰兹对偶群作用在 $\ell$-进层范畴上。}{I：本工程 glossary.tsv；A：原文 §I.2}
：它们将同时使 Weil 群 $W_E$（即 $\pi_1(\Spec\breve{E}/\varphi^\ZZ)$）以及（经过细致的研究）朗兰兹对偶群 $\widehat{G}$ 出现。回忆 Hecke 算子与 $G$-主丛的修正有关，并由修正发生处的曲线上的点 $x$ 以及修正在该点处的类型来参数化（后者可以用 $G$ 的余特征以组合的方式编码——这最终导致 $\widehat{G}$ 的出现）。通常，当 $x$ 变动时，Hecke 算子的效应是局部常的。在这种情形下，让 $x$ 变动就等于让 $\pi_1(X)$ 作用，其中 $X$ 是相应的曲线；因此这条曲线现在应当是 $\Spec\breve{E}/\varphi^\ZZ$。

于是，若要定义 Hecke 算子，我们需要能够考虑 $G$-等晶体的修正。这些修正理应发生在 $\Spec R((t))\to\Spec R$ 的一个截面处（resp.\ 一个并不存在的映射 $\Spec(W(R)\otimes_{W(\FF_q)}E)\to\Spec R$ 的截面处）。不幸的是，映射 $R\to R((t))$ 并\srcpageinl{12}不承认任何截面。事实上，我们当然想考虑连续截面；这样的连续截面于是将与映射 $\ol{\FF}_q((t))=\breve{E}\to R$ 一一对应。换言之，与上一段的动机一致，相应的曲线应当是 $\Spec\breve{E}$，或者更确切地说是 $\Spec\breve{E}$ 模 Frobenius——于是我们可以自然地期望按上述配方得到 $\pi_1(\Spec\breve{E}/\varphi^\ZZ)$ 的作用。

然而，为了使这幅图景得以实现，我们需要处在有连续映射 $\FF_q((t))\to R$ 的情形。换言之，只有把自身置于 $R$ 本身就是某种 Banach 环的设定中，我们才能指望有截面。

这最终把我们带到本文所考虑的设定。即：我们把完满 $\FF_q$-概形的范畴换成 $\FF_q$ 上完满空间的范畴 $\mathrm{Perf}=\mathrm{Perf}_{\FF_q}$\yizhu{perfectoid space 本书统一译作「完满空间」（perfect 的「完满」＋ -oid 后缀），不作音译；「完满域/完满环」同此。}{I：本工程 glossary.tsv；A：原文 §I.2}
。局部地，任一 $S\in\mathrm{Perf}$ 都有形式 $S=\Spa(R,R^+)$，其中 $R$ 是完满 Tate $\FF_q$-代数：意思是 $R$ 是一个完满拓扑代数，存在拓扑幂零单位 $\varpi\in R$（称为伪一致化子）使它成为 $\FF_q((\varpi))$ 上的 Banach 代数。此外 $R^+\subset R$ 是开的、整闭的幂有界元子环。通常 $R^+=R^\circ$ 就是幂有界元构成的子环，我们因而采用缩写 $\Spa R=\Spa(R,R^\circ)$。$S$ 的几何（秩 $1$）点由 $\Spa C$ 给出，其中 $C$ 是完备代数闭非阿基米德域；一如往常，理解几何点是关键的第一步。关于 adic 空间与完满空间的导引，我们参见 [SW20]。

对任意 $S=\Spa(R,R^+)$，我们需要定义 $\Spec R((t))/\varphi^\ZZ$ 的类比，并顾及 $R$ 的拓扑。注意对离散的 $R'$，有
\[
  \Spa R'((t)) = \Spa R'\times_{\Spa\FF_q}\Spa\FF_q((t)),
\]
而我们总可以自由地把 $\Spec R'((t))/\varphi^\ZZ$ 换成 $\Spa R'((t))/\varphi^\ZZ$，因为二者有相同的向量丛范畴。这就提示 $\Spec R'((t))$ 的类比是
\[
  \Spa(R,R^+)\times_{\Spa\FF_q}\Spa\FF_q((t)) = \mathbb{D}^*_{\Spa(R,R^+)},
\]
即 $\Spa(R,R^+)$ 上的一个去心开单位圆盘，坐标为 $t$。注意
\[
  \Spa(R,R^+)\times_{\Spa\FF_q}\Spa\FF_q((t)) \subset \Spa R^+\times_{\Spa\FF_q}\Spa\FF_q[[t]] = \Spa R^+[[t]]
\]
是 $t$ 与 $\varpi\in R^+$ 可逆的轨迹，其中 $\varpi$ 是 $R$ 的一个拓扑幂零单位。后一定义可以推广到混合特征：我们令
\[
  \Spa(R,R^+)\dot\times_{\Spa\FF_q}\Spa E \subset \Spa R^+\dot\times_{\Spa\FF_q}\Spa O_E := \Spa W_{O_E}(R^+)
\]
为使 $\pi$ 与 $[\varpi]\in W_{O_E}(R^+)$ 可逆的开子集。此空间与 $\varpi$ 的选取无关，因为对任何另一个这样的 $\varpi'$，存在 $n\ge1$ 使 $\varpi\mid\varpi'^n$ 且 $\varpi'\mid\varpi^n$，而它们的 Teichmüller 代表也如此。我们要指出，符号 $\dot\times$ 纯属记号：当然不存在 adic 空间之间的映射 $\Spa E\to\Spa\FF_q$，使纤维积得以取到。

\begin{definition}% 原编号 I.2.2\label{def:I.2.2}
Fargues–Fontaine 曲线（对局部域 $E$ 与 $S=\Spa(R,R^+)\in\mathrm{Perf}$）是 $E$ 上如下定义的 adic 空间：
\[
  X_S = X_{S,E} = \bigl(\Spa(R,R^+)\times_{\Spa\FF_q}\Spa\FF_q((t))\bigr)/\varphi^\ZZ,
\]
相应地，
\[
  X_S = X_{S,E} = \bigl(\Spa(R,R^+)\dot\times_{\Spa\FF_q}\Spa E\bigr)/\varphi^\ZZ.
\]
\leavevmode\srcpage{13}其中 Frobenius $\varphi$ 作用于 $(R,R^+)$。
\end{definition}

与 $G$-等晶体的讨论相比，一个新的特点是：$\varphi$ 的作用是自由且完全间断的，因此按 $\varphi$ 取商在 adic 空间的范畴中是有定义的。事实上，在
\[
  Y_S = \Spa(R,R^+)\dot\times_{\Spa\FF_q}\Spa E \subset \Spa W_{O_E}(R^+)
\]
上可以比较 $\pi$ 与 $[\varpi]$ 的绝对值。由于二者都是拓扑幂零单位，比值
\[
  \mathrm{rad} = \log(|[\varpi]|)/\log(|\pi|) : |Y_S| \to (0,\infty)
\]
给出一个良定义的连续映射。$|Y_S|$ 上的 Frobenius 把 $\mathrm{rad}$ 乘以 $q$，这证明了该作用是自由且完全间断的。

我们指出，在函数域情形 $E=\FF_q((t))$，空间
\[
  Y_S = S\times_{\Spa\FF_q}\Spa\FF_q((t)) = \mathbb{D}^*_S
\]
恰是 $S$ 上的一个去心开单位圆盘。在这幅图景中，半径函数量度的是到原点的距离：靠近原点时，半径映射接近 $0$；而靠近开单位圆盘的边界时，它接近 $\infty$。然而按 $\varphi$ 取商不再是 $S$ 上的 adic 空间，因为 $\varphi$ 作用在 $S$ 上。于是
\[
  X_S = Y_S/\varphi^\ZZ = \mathbb{D}^*_S/\varphi^\ZZ
\]
局部地是 $S$ 上有限型的 adic 空间，但整体上不是。当 $S=\Spa C$ 为几何点时，此空间最先由 Hartl–Pink [HP04] 研究。

若 $S=\Spa C$ 是几何点而 $E$ 是一般的，则这条曲线（或者更确切地说，一个与之密切相关的概形版本）已由 Fargues–Fontaine [FF18] 广泛研究，其中证明了它在 $p$-进 Hodge 理论中起核心作用。从 adic 空间的角度看，它已由 Kedlaya–Liu [KL15] 研究。特别地，在 $S$ 是一个点的这种情形下，$X_S$ 确为一条曲线：它是一个强 Noether 的 adic 空间，其连通的 affinoid 子集是主理想整环的谱。特别地，在这种情形下有一个性质良好的「经典点」概念\yizhu{「经典点」（classical point）指局部对应于环的极大理想的点；对曲线而言它们相当于通常的闭点，但因 $X_S$ 整体只是 adic 空间（未必强 Noether），须逐点局部定义。}{A：原文 §I.3 与 pp.51, 56；I：本工程 glossary.tsv}
，指的是局部对应于极大理想的那些点。这些点可以被分类。在等特征情形，
\[
  Y_S = \mathbb{D}^*_S = S\times_{\Spa\FF_q}\Spa\FF_q((t))
\]
的描述表明，闭点与映射 $S\to\Spa\FF_q((t))$ 模 Frobenius 一一对应；只不过这里还须按 $t\mapsto t^q$ 取商。在混合特征情形，情况更为微妙，并引向完满空间的倾斜构造。

\begin{proposition}% 原编号 I.2.3\label{prop:I.2.3}
若 $E$ 为混合特征且 $S=\Spa C$ 是几何点，则 $X_C$ 的经典点与 $C$ 在 $E$ 上的 untilt $C^\sharp$ 模 Frobenius 作用一一对应。
\end{proposition}

这里我们回忆，对任何完备代数闭域 $C'|E$，更一般地对任何完满 Tate 环 $R$，都可以构造倾斜
\[
  R^\flat = \varprojlim_{x\mapsto x^p} R,
\]
其中加法在整元环上通过双射 $R^{\flat+}=\varprojlim_{x\mapsto x^p}R^+\cong\varprojlim_{x\mapsto x^p}R^+/\pi$ 来定义，此时 $x\mapsto x^p$ 与 $R^+/\pi$ 上的加法相容。于是 $R^\flat$ 是特征 $p$ 的完满 Tate 代数。几何上，把 $\Spa(R,R^+)$ 送到 $\Spa(R^\flat,R^{\flat+})$\srcpageinl{14}定义了完满空间上的倾斜函子 $T\mapsto T^\flat$，它保持底层的拓扑空间与 \text{ét}ale 站点，参看 [SW20]。

可以看到，当 $S=\Spa C$ 为几何点时，$X_S$ 的经典点与 $S$ 的 untilt $S^\sharp$ 连同映射 $S\to\Spa E$ 模 Frobenius 作用一一对应。回忆 [SW20]：对 $W(\FF_q)$ 上任何 adic 空间 $Z$，定义函子
\[
  Z^\diamond : \mathrm{Perf}\to\mathrm{Sets} : S\mapsto\{S^\sharp,\, f : S^\sharp\to Z\},
\]
它把 $\ol{\FF}_q$ 上完满空间 $S$ 送到一对对象，即 $S$ 的一个 untilt $S^\sharp$ 以及一个映射 $S^\sharp\to Z$。若 $Z$ 是解析 adic 空间，则 $Z^\diamond$ 是一个菱形\yizhu{「菱形」（diamond）指完满空间按 pro-étale 等价关系取的商；它是 $v$-层中「表现得像完满空间」的那一类，经典点是其重要特例。}{A：原文 §I.2 及 [SW20] 的定义；I：本工程 glossary.tsv}
，即一个完满空间按 pro-\text{ét}ale 等价关系取的商。于是 $X_S$ 的经典点与菱形
\[
  (\Spa\breve{E})^\diamond/\varphi^\ZZ
\]
的 $S$-值点一一对应。
更一般地，对任何 $S\in\mathrm{Perf}$，映射 $S\to(\Spa\breve{E})^\diamond/\varphi^\ZZ$ 与次数 $1$ 的 Cartier 除子 $D_S\subset X_S$ 一一对应，因此我们定义
\[
  \Div^1 = (\Spa\breve{E})^\diamond/\varphi^\ZZ.
\]
我们要提醒读者，此处 Frobenius 的作用是几何 Frobenius。特别地，当 $E$ 为混合特征时，它只存在于 $(\Spa\breve{E})^\diamond$ 上，而不存在于 $\Spa\breve{E}$ 上。然而仍有
\[
  \pi_1(\Div^1) = W_E.
\]

至此，我们对 $\Spec E$ 的几何的这一长串思考结束了：我们抵达了 Fargues–Fontaine 曲线，以及它的各种化身。为便于读者定位，我们在此把它们回顾一遍：

\noindent (i) 对任何完备代数闭非阿基米德域 $C|\FF_q$，有曲线 $X_C=X_{C,E}$，它是 $E$ 上强 Noether 的 adic 空间，局部地是主理想整环的 adic 谱。人们也可以构造一个概形版本 $X_C^{\mathrm{alg}}$，它有相同的经典点与相同的向量丛范畴。经典点与 $C$ 的 untilt $C^\sharp|E$ 模 Frobenius 一一对应。

\noindent (ii) 更一般地，对任何完满空间 $S\in\mathrm{Perf}$，有「曲线族」$X_S$，它仍是 $E$ 上的 adic 空间，但不再是强 Noether 的。若 $S$ 是 affinoid 的，则有概形版本 $X_S^{\mathrm{alg}}$，具有相同的向量丛范畴。

\noindent (iii)「镜像曲线」$\Div^1=(\Spa\breve{E})^\diamond/\varphi^\ZZ$，它只是一个菱形。对任何 $S\in\mathrm{Perf}$，它参数化「$X_S$ 上次数 $1$ 的 Cartier 除子」。

这里有一个奇特的现象：并不存在一条「绝对曲线」，使其余各条都是它的基变换。另一个奇特之处是：次数 $1$ 的 Cartier 除子所构成的空间并不是曲线本身。

又到了研究 $G$-主丛的时候。这引出下面的定义。

\begin{definition}% 原编号 I.2.4\label{def:I.2.4}
令
\[
  \Bun_G : \mathrm{Perf}\to\{\text{groupoids}\} : S\mapsto\{G\text{-torsors on }X_S\}
\]
为 Fargues–Fontaine 曲线上 $G$-主丛的模叠。
\end{definition}

\begin{remark}% 原编号 I.2.5\label{rem:I.2.5}
我们在此强调：虽然「Fargues–Fontaine 曲线」并不是一个真正良定义的概念，但「Fargues–Fontaine 曲线上 $G$-主丛的模叠」是。
\end{remark}

\leavevmode\srcpage{15}由于 $X_S$ 映向 $\Spa\breve{E}/\varphi^\ZZ$，存在自然的拉回函子 $G$-$\Isoc\to\Bun_G(S)$。下面的结果在多数情形下归功于 Fargues [Far20a]，并由 Anschütz [Ans19] 补全。

\begin{theorem}% 原编号 I.2.6\label{thm:I.2.6}
若 $S=\Spa C$ 是几何点，则映射
\[
  B(G)\to\Bun_G(S)/\!\cong
\]
是双射。特别地，$X_S$ 上任何向量丛都是与 $D_{-\lambda}$ 相联系的向量丛 $O_{X_S}(\lambda)$ 的直和，其中 $\lambda\in\QQ$。

在此双射下，$b\in B(G)$ 是基本的，当且仅当 $X_S$ 上相应的 $G$-主丛 $\mathcal{E}_b$ 在 Atiyah–Bott [AB83] 意义下是半稳定的。
\end{theorem}

然而，自同构群不再相同。在叠的层面上，我们有下面的结果。

\begin{theorem}% 原编号 I.2.7\label{thm:I.2.7}
prestack $\Bun_G$ 是 v-叠。它容许一个由局部闭子叠构成的分层
\[
  i^b : \Bun_G^b \subset \Bun_G
\]
（$b\in B(G)$），由在每个几何点处 $G$-主丛同构于 $\mathcal{E}_b$ 的轨迹组成。此外，每个层
\[
  \Bun_G^b \cong [\ast/\widetilde{G}_b]
\]
都是群 $\widetilde{G}_b$ 的分类空间，而 $\widetilde{G}_b$ 是局部 profinite 群 $G_b(E)$ 被一个「幂幺群菱形」\yizhu{「unipotent group diamond」（幂幺群菱形）在本篇只作描述性使用并加引号，原文未给定义；本译本查遍全书未见它的正式定义，暂按字面译出，此处存疑（见 `errata.tsv` E6）。}{I：全库检索未见定义；A：原文 p.15 首次出现、p.90 只用到同类概念}
的扩张。

半稳定轨迹 $\Bun_G^{\mathrm{ss}}\subset\Bun_G$ 是开子叠，且
\[
  \Bun_G^{\mathrm{ss}} = \bigsqcup_{b\in B(G)_{\mathrm{basic}}}[\ast/G_b(E)].
\]
\end{theorem}

\begin{remark}% 原编号 I.2.8\label{rem:I.2.8}
此定理在形式上与定理 I.2.1 极为相似。然而有一个关键差别，即闭包关系被反转了：对 $\Bun_G$ 而言，当 $b\in B(G)$ 基本时，$\Bun_G^b$ 的包含是开浸入；而在定理 I.2.1 中它是闭浸入。请注意，基本的 $b\in B(G)$ 对应于半稳定的 $G$-丛，而人们确实会期望半稳定轨迹是开的。一般地，$\Bun_G$ 的表现很像射影直线上 $G$-丛的叠。
\end{remark}

\begin{remark}% 原编号 I.2.9\label{rem:I.2.9}
我们在完满设定下定义了 Artin 叠的概念，而 $\Bun_G$ 确实是 Artin 叠；关于 $\Bun_G$ 上几何结果的更详细描述，我们参见 §I.4。这表明 $\Bun_G$ 的有限性远远好于 $G$-$\Isoc$，即便它定义在更 exotic 的测试对象上。
\end{remark}

我们可以定义 $\Bun_G$ 上 $\ell$-进层的导出范畴
\[
  D(\Bun_G,\ol{\QQ}_\ell).
\]
它容许一个半正交分解，其各项为所有 $D(\Bun_G^b,\ol{\QQ}_\ell)$，并且
\[
  D(\Bun_G^b,\ol{\QQ}_\ell)\cong D([\ast/G_b(E)],\ol{\QQ}_\ell)\cong D(G_b(E),\ol{\QQ}_\ell)
\]
等价于 $G_b(E)$ 的光滑表示的导出范畴。

\begin{remark}% 原编号 I.2.10\label{rem:I.2.10}
可以合理地期望这个范畴等价于 Xiao–Zhu 所定义的范畴 $D(G\text{-}\mathcal{I}\mathrm{soc},\ol{\QQ}_\ell)$。然而本文不追求这一比较。
\end{remark}

\leavevmode\srcpageinl{16}最后，我们可以定义那个把所有关键角色汇聚到一起的 Hecke 叠。考虑全局 Hecke 叠 $\Hck_G$，它参数化 $X_S$ 上 $G$-丛的对 $(\mathcal{E},\mathcal{E}')$，连同给出 $X_S$ 中一个次数 $1$ 的 Cartier 除子 $D_S\subset X_S$ 的映射 $S\to\Div^1$，以及一个同构
\[
  f : \mathcal{E}|_{X_S\setminus D_S}\cong\mathcal{E}'|_{X_S\setminus D_S},
\]
该同构沿 $D_S$ 是亚纯的。这给出一个对应
\[
  \Bun_G \xleftarrow{\ h_1\ } \Hck_G \xrightarrow{\ h_2\ } \Bun_G\times\Div^1.
\]
为定义 Hecke 算子，我们需要界住修正，即界住 $f$ 沿 $D_S$ 的极点。这由局部 Hecke 叠 $\ul{\Hck}_G$ 描述，它参数化 $X_S$ 沿 $D_S$ 的完备化上的 $G$-主丛对，连同一个离开 $D_S$ 的同构；于是存在从全局 Hecke 叠到局部 Hecke 叠的自然映射 $\ul{\Hck}_G\to\Hck_G$。几何上，$\ul{\Hck}_G$ 容许一个按 $G$ 的余特征的共轭类给出的 Schubert 分层；特别地，对每个共轭类 $\mu:\mathbb{G}_m\to G$，有闭 Schubert 胞腔 $\ul{\Hck}_{G,\le\mu}$。通过拉回，这定义了一个对应
\[
  \Bun_G \xleftarrow{\ h_{1,\le\mu}\ } \ul{\Hck}_{G,\le\mu} \xrightarrow{\ h_{2,\le\mu}\ } \Bun_G\times\Div^1,
\]
此时 $h_{1,\le\mu}$ 与 $h_{2,\le\mu}$ 都是紧合的。于是可以考虑 Hecke 算子
\[
  Rh_{2,\le\mu,\ast}h_{1,\le\mu}^\ast : D(\Bun_G,\Lambda)\to D(\Bun_G\times\Div^1,\Lambda).
\]

下面的定理保证：当 $\Div^1$ 上的点变动时，Hecke 算子必然是局部常的，从而给出 $\pi_1(\Div^1)=W_E$ 的表示。在下文中，我们对所采用的 $D(-,\ol{\QQ}_\ell)$ 的精确定义以及 $W_E$-等变对象的概念有些随意：细节上的小心处理在正文中给出。

\begin{theorem}% 原编号 I.2.11\label{thm:I.2.11}
沿映射 $\Div^1\to[\ast/W_E]$ 的拉回诱导出等价
\[
  D(\Bun_G\times\Div^1,\ol{\QQ}_\ell)\cong D(\Bun_G\times[\ast/W_E],\ol{\QQ}_\ell)\cong D(\Bun_G,\ol{\QQ}_\ell)^{BW_E}.
\]
\end{theorem}

于是，Hecke 算子产出 $D(\Bun_G,\ol{\QQ}_\ell)$ 中的 $W_E$-等变对象，使 Weil 群自然地出现。

人们还想理解 Hecke 算子如何复合。这自然引向把 $D(\Hck_G,\ol{\QQ}_\ell)$ 作为卷积下的幺半范畴来研究。这里我们有几何 Satake 等价\yizhu{「几何 Satake 等价」把朗兰兹对偶群 $\widehat{G}$ 的表示范畴与 Hecke 叠上（反常）层的卷积幺半范畴等同起来；它是几何 Langlands 纲领的支柱之一。}{A：原文 §I.6 与 [MV07]；I：本工程 glossary.tsv}
。在通常的光滑射影曲线（定义在 $\CC$ 上）的设定中，这由 Lusztig [Lus83]、Ginzburg [Gin90] 与 Mirković–Vilonen [MV07] 的论文建立。下面的定理只是一个初步的近似；我们实际上将证明一个更精确的版本，它带 $\ZZ_\ell$-系数，描述 $\Hck_G$ 上所有的反常层，并适用于 Beilinson–Drinfeld Grassmann 簇，其精神一如 Gaitsgory 的论文 [Gai07]。

\begin{theorem}% 原编号 I.2.12\label{thm:I.2.12}
存在从 $\Rep\widehat{G}$ 到 $D(\Hck_G,\ol{\QQ}_\ell)$ 的自然幺半函子。
\end{theorem}

\begin{remark}% 原编号 I.2.13\label{rem:I.2.13}
我们对定理 I.2.12 的证明遵循 Mirković–Vilonen 证明的策略，特别地，借助融合积在反常层的范畴上定义了一个自然的对称幺半结构。这要求人们在基曲线的若干份拷贝上工作，并\srcpageinl{17}让各个点相互碰撞。先验地，这样的做法在混合特征中可行是非常令人惊讶的，因为它需要一个像 $\Spa\QQ_p\dot\times_{\Spa\FF_p}\Spa\QQ_p$ 这样的空间。然而这类空间作为菱形是存在的，而这是 [SW20] 的主要创新之一。
\end{remark}

\begin{remark}% 原编号 I.2.14\label{rem:I.2.14}
利用局部 Hecke 叠（它本质上就是 [SW20] 的 $B^+_{\mathrm{dR}}$-仿射 Grassmann 簇）到 Witt 向量仿射 Grassmann 簇的退化，定理 I.2.12 给出 Zhu 关于 Witt 向量仿射 Grassmann 簇的几何 Satake 等价 [Zhu17] 的一个新证明。事实上，我们甚至证明了一个带 $\ZZ_\ell$-系数的版本，从而也重新得到了 Yu [Yu22] 的结果。
\end{remark}

\begin{remark}% 原编号 I.2.15\label{rem:I.2.15}
关于 $\ell$-进层的形式化，我们要提醒读者：在定理 I.2.12 的表述中我们略有作弊；上文隐含的 $D(\Bun_G,\ol{\QQ}_\ell)$ 的定义与定理 I.2.12 中所隐含的并不是同一个。若改用挠系数，问题就会消失；况且无论如何，这些问题本质上都是技术性的。
\end{remark}

于是，这也使朗兰兹对偶群 $\widehat{G}$ 自然地出现。对 $\widehat{G}$ 的任何表示 $V$，我们得到一个 Hecke 算子
\[
  T_V : D(\Bun_G,\ol{\QQ}_\ell)\to D(\Bun_G,\ol{\QQ}_\ell)^{BW_E}.
\]
此外，Hecke 算子彼此交换，并且
\[
  T_{V\otimes W}\cong T_V\circ T_W|_{\Delta(W_E)},
\]
这里我们注意到 $T_V\circ T_W$ 自然地取值于 $W_E\times W_E$-等变对象；右端的限制指的是限制到对角拷贝 $\Delta(W_E)\subset W_E\times W_E$ 的作用上。

至此，$G(E)$ 的表示论（它全忠实地落在 $D(\Bun_G,\ol{\QQ}_\ell)$ 之中）、Weil 群 $W_E$ 与对偶群 $\widehat{G}$ 三者自然地相互作用。事实证明，正是这一范畴结构足以对 $D(\Bun_G,\ol{\QQ}_\ell)$ 中的（Schur-）不可约对象 $A$——特别是对 $G(E)$ 的不可约光滑表示——构造 $L$-参数。我们将在下文 §I.9 讨论 $L$-参数的构造。

我们指出，整幅图景与 Nadler–Yun [NY19] 所考虑的 Betti 几何 Langlands 情形完全平行；事实上，整个策略可以描述为「Fargues–Fontaine 曲线上的几何 Langlands 纲领」。颇为有趣的是：我们的追求原本是在算术设定中理解局部 Langlands 对应，且是对可能深度分歧的表示；最终我们却是通过把它联系到几何设定中、处处无分歧的整体 Langlands 对应来解决它的。

在本导引的余下部分，我们对这幅图景的各个方面给出更详细的概述：

\noindent (i) Fargues–Fontaine 曲线（§I.3）；

\noindent (ii) 叠 $\Bun_G$ 的几何（§I.4）；

\noindent (iii) $\Bun_G$ 上 $\ell$-进层的导出范畴（§I.5）；

\noindent (iv) 几何 Satake 等价（§I.6）；

\noindent (v) Rapoport–Zink 空间、局部 Shimura 簇以及更一般的 shtuka 模空间的上同调的有限性（§I.7）；

\noindent (vi) $L$-参数叠（§I.8）；

\noindent (vii) $L$-参数的构造（§I.9）；

\noindent\srcpageinl{18}(viii) 谱作用（§I.10）；

\noindent (ix) 本文所阐述的思想的来源（§I.11）。

上述各项大体上对应本文的各章，而每一章开头都有对这些导引的重述。

\section{Fargues–Fontaine 曲线}% 原文 I.3. The Fargues–Fontaine curve

Fargues–Fontaine 曲线已在 Fargues–Fontaine 的书 [FF18] 中被广泛研究，而更多的结果——特别是在相对情形下——已由 Kedlaya–Liu [KL15] 得到。在第一章中，我们重新证明这些基础性的结果，从而也收集并统一了一些（往往只对 $E=\QQ_p$ 证明过的）结果。

首批结果涉及当 $S=\Spa C$（$C|\FF_q$ 为某个完备代数闭非阿基米德域）时的 Fargues–Fontaine 曲线 $X_C=X_S$。在这种情形下我们定义 $X_C$ 的经典点的概念；它们构成 $|X_C|$ 的一个子集。$X_C$ 的基本有限性性质概括在下面的结果中。

\begin{theorem}% 原编号 I.3.1\label{thm:I.3.1}
adic 空间 $X_C$ 局部地是 adic 谱 $\Spa(B,B^+)$，其中 $B$ 是主理想整环；$\Spa(B,B^+)\subset X_C$ 的经典点与 $B$ 的极大理想一一对应。对每个经典点 $x\in X_C$，$x$ 的剩余域是 $C$ 在 $E$ 上的一个 untilt $C^\sharp$，这诱导出 $X_C$ 的经典点与 $C$ 在 $E$ 上的 untilt $C^\sharp$（模 Frobenius 作用）之间的一一对应。
\end{theorem}

在等特征情形，定理 I.3.1 是表述 $X_C=\mathbb{D}^*_C/\varphi^\ZZ$ 与刚性解析几何中经典结果的直接推论。在 $p$-进情形，我们利用倾斜把它归约到等特征情形。在一个关键的转折处，为理解 $X_C$ 的 Zariski 闭子集，我们使用了「Zariski 闭蕴含强 Zariski 闭」这一结果 [BS22]。借助这些想法，我们得以给出一个基本上无需计算的证明。

一个关键结果是向量丛的分类。

\begin{theorem}% 原编号 I.3.2\label{thm:I.3.2}
从 $\Isoc_E$ 到 $X_C$ 上向量丛的函子在同构类上诱导出双射。特别地，对任何斜率 $\lambda\in\QQ$，存在唯一的稳定向量丛 $O_{X_C}(\lambda)$，而任何向量丛 $\mathcal{E}$ 都可以写成稳定向量丛的直和。
\end{theorem}

我们给出定理 I.3.2 的一个新的、自足的证明，其中关键地使用了 [Sch17a] 与 [SW20] 中得到的向量丛的 v-下降结果，以及本文所建立的关于 Banach–Colmez 空间几何的基本结果。Hartl–Pink [HP04] 在等特征情形的证明与 Kedlaya 在 $p$-进情形的证明 [Ked04] 依赖于繁重的计算，而 Fargues–Fontaine [FF18] 的证明则依赖于对 $\pi$-可除 $O$-模的 Lubin–Tate 模空间与 Drinfeld 模空间的描述。我们的证明与 Colmez 在 [Col02] 中的论证有关。

允许一般的 $S\in\mathrm{Perf}_{\FF_q}$，我们把次数 $1$ 的 Cartier 除子的模空间定义为 $\Div^1=\mathrm{Spd}\,\breve{E}/\varphi^\ZZ$。给定映射 $S\to\Div^1$，可以定义与之相伴的闭 Cartier 除子 $D_S\subset X_S$；局部地，它由 $S$ 在 $E$ 上的一个 untilt $D_S=S^\sharp\subset X_S$ 给出，而这把 $\Div^1$ 嵌入到 $X_S$ 上闭 Cartier 除子的空间之中（这为其名称提供了依据）。另一个重要的结果是下面的丰沛性结果，参看 [KL15, Proposition 6.2.4]，它蕴含人们可以定义曲线的代数版本，且该版本容许同样的向量丛理论。

\begin{theorem}% 原编号 I.3.3\label{thm:I.3.3}
\leavevmode\srcpage{19}设 $S\in\mathrm{Perf}$ 是 affinoid 的。对 $X_S$ 上任何向量丛 $\mathcal{E}$，当 $n\gg0$ 时扭 $\mathcal{E}(n)$ 由整体截面生成，且没有高阶上同调。定义分次环
\[
  P = \bigoplus_{n\ge0}H^0(X_S,\cO_{X_S}(n))
\]
与概形 $X_S^{\mathrm{alg}}=\mathrm{Proj}\,P$，则存在局部环化空间的自然映射 $X_S\to X_S^{\mathrm{alg}}$，沿它拉回定义了向量丛范畴的一个等价，并保持上同调。

若 $S=\Spa C$，其中 $C$ 是某个完备代数闭非阿基米德域，则 $X_C^{\mathrm{alg}}$ 是 Krull 维数为 $1$ 的正则 Noether 概形，局部地是主理想整环的谱，而其闭点与 $X_C$ 的经典点一一对应。
\end{theorem}

我们还需要理解向量丛的族，即一般 $S$ 上 $X_S$ 的向量丛 $\mathcal{E}$。这里的主要结果是下面这个。

\begin{theorem}% 原编号 I.3.4\label{thm:I.3.4}
设 $S\in\mathrm{Perf}$，$\mathcal{E}$ 为 $X_S$ 上的向量丛。则把点 $s\in S$ 送到 $\mathcal{E}|_{X_s}$ 的 Harder–Narasimhan 多边形的函数定义了 $S$ 上的一个半连续函数。若该函数为常值，则 $\mathcal{E}$ 容许一个整体的 Harder–Narasimhan 分层，且在 $S$ 上 pro-\text{ét}ale 局部地可以找到与 $\bigoplus\cO_{X_S}(\lambda)$ 的同构。

特别地，若 $\mathcal{E}$ 处处是斜率 $0$ 的半稳定向量丛，则 $\mathcal{E}$ 是 pro-\text{ét}ale 局部平凡的，而这类 $\mathcal{E}$ 的范畴等价于 $S$ 上 pro-\text{ét}ale $\underline{E}$-局部系统的范畴。
\end{theorem}

证明定理 I.3.4 的关键在于构造 $\mathcal{E}$ 的某些整体截面。为此，我们使用 v-下降技术，并对 $\mathcal{E}$ 的整体截面所构成的空间进行分析；这些空间被称为 Banach–Colmez 空间，最初由 [Col02] 引入（术语略有不同）。

\begin{definition}% 原编号 I.3.5\label{def:I.3.5}
设 $\mathcal{E}$ 为 $X_S$ 上的向量丛。与 $\mathcal{E}$ 相伴的 Banach–Colmez 空间 $\mathrm{BC}(\mathcal{E})$ 是 $S$ 上的一个局部空间菱形，对 $T\in\mathrm{Perf}_S$，其 $T$-值点由
\[
  \mathrm{BC}(\mathcal{E})(T) = H^0(X_T,\mathcal{E}|_{X_T})
\]
给出。类似地，若 $\mathcal{E}$ 处处只有负的 Harder–Narasimhan 斜率，则负 Banach–Colmez 空间 $\mathrm{BC}(\mathcal{E}[1])$ 是 $S$ 上的一个局部空间菱形，其 $T$-值点为
\[
  \mathrm{BC}(\mathcal{E}[1])(T) = H^1(X_T,\mathcal{E}|_{X_T}).
\]
\end{definition}

这里隐含的是：该函子确实定义了一个局部空间菱形。为此我们计算 Banach–Colmez 空间的一些关键例子。例如，若 $\mathcal{E}=\cO_{X_S}(\lambda)$ 且 $0<\lambda\le[E:\QQ_p]$（resp.\ 若 $E$ 为等特征，则对一切正的 $\lambda$），则 $\mathrm{BC}(\mathcal{E})$ 可由一个完满开单位圆盘表示（其维数由 $\lambda$ 的分子给出）。一个特殊情形是 $\mathrm{BC}(\cO_{X_S}(1))$ 与某个 Lubin–Tate 形式群律的万有覆盖的等同，这给出 Lubin–Tate 理论（因而局部类域论）与 Fargues–Fontaine 曲线之间的密切关系；另见 [Far20b]。另一方面，对更大的 $\lambda$ 或负的 $\lambda$，Banach–Colmez 空间是更 exotic 的对象；例如负 Banach–Colmez 空间
\[
  \mathrm{BC}(\cO_{X_C}(-1)[1])\cong(\mathbb{A}^1_{C^\sharp})^\diamond/E
\]
\srcpageinl{20}是按 $\ul{E}\subset\mathbb{A}^1_{C^\sharp}$ 的平移作用对仿射线取的商。我们要指出，我们对分类定理（定理 I.3.2）的证明最终依赖于一个负面结果，即 $\mathrm{BC}(\cO_{X_C}(-1)[1])$ 不能由完满空间表示！\footnote{事实上，我们只在 $E$ 为 $p$-进时确知这一点；在函数域情形，我们补充一个小论证来绕开这个问题。}

对定理 I.3.4 的证明而言，一个关键结果是：射影化的 Banach–Colmez 空间
\[
  (\mathrm{BC}(\mathcal{E})\setminus\{0\})/E^\times
\]
是紧合的——它们是「$S$ 上射影空间族」的适当类比。特别地，它们在 $S$ 中的像是闭子集；而若该像就是整个 $S$，则在取一个 v-覆盖之后可以找到 $\mathcal{E}$ 的一个无处为零的截面，因为此时射影化的 Banach–Colmez 空间是 $S$ 的 v-覆盖。由此，定理 I.3.4 就容易得到了。

\section{$\Bun_G$ 的几何}% 原文 I.4. The geometry of Bun_G

我们来讨论 $\Bun_G$ 的几何。这里 $G$ 可以是任何非阿基米德局部域 $E$ 上的约化群，其剩余域 $\FF_q$ 的特征为 $p$。回忆 Kottwitz 的 $G$-等晶体集合 $B(G)=B(E,G)$，它可以由两个离散不变量以组合的方式描述。第一个是 Newton 点
\[
  \nu : B(G)\to(X_*(T)^+_\QQ)^\Gamma,
\]
其中 $T$ 是 $G$ 的泛 Cartan，$\Gamma=\Gal(\ol{E}|E)$。更确切地说，任何 $G$-等晶体 $\mathcal{E}$ 都定义一个斜率态射 $\mathbb{D}\to G_{\breve{E}}$，其中 $\mathbb{D}$ 是以 $\QQ$ 为余特征群的对角化群；其定义归约到 $\mathrm{GL}_n$ 的情形，在那里它相当于等晶体的斜率分解。$G$-等晶体的同构给出共轭的斜率态射，由此定义了映射 $\nu$。

另一个映射是 Kottwitz 不变量
\[
  \kappa : B(G)\to\pi_1(G_{\ol{E}})_\Gamma.
\]
它的定义是间接的：从环面出发，过渡到导群单连通的 $G$ 的情形，最后用 z-扩张过渡到一般情形。于是 Kottwitz 证明了
\[
  (\nu,\kappa) : B(G)\to(X_*(T)^+_\QQ)^\Gamma\times\pi_1(G_{\ol{E}})_\Gamma
\]
是单射。此外，$\kappa$ 诱导出 $B(G)_{\mathrm{basic}}$ 与 $\pi_1(G_{\ol{E}})_\Gamma$ 之间的双射\yizhu{此处 $\kappa$ 是经 z-扩张间接定义的；其与 $\pi_1$ 的完整对应（含基本集上为双射）在 §III.2 中正式建立。}{A：原文 §III.2；I：本工程 glossary.tsv}
。非基本的元素可以用 Levi 子群来描述。

利用 $\nu$ 与 $\kappa$，可以按如下方式在 $B(G)$ 上定义偏序：若 $\kappa(b)=\kappa(b')$ 且 $\nu_b\le\nu_{b'}$（关于支配序），则称 $b\le b'$。

差一个符号，人们可以把 $\nu$（resp.\ $\kappa$）看作一个 $G$-丛的 Harder–Narasimhan 多边形（resp.\ 第一 Chern 类）。

\begin{theorem}% 原编号 I.4.1\label{thm:I.4.1}
prestack $\Bun_G$ 满足以下性质。

\noindent (i) prestack $\Bun_G$ 对 v-拓扑是叠。

\noindent (ii) 点 $|\Bun_G|$ 自然地与 Kottwitz 的 $G$-等晶体集合 $B(G)$ 一一对应。

\noindent\srcpage{21}(iii) 映射
\[
  \nu : |\Bun_G|\to B(G)\to(X_*(T)^+_\QQ)^\Gamma
\]
是半连续的，而
\[
  \kappa : |\Bun_G|\to B(G)\to\pi_1(G_{\ol{E}})_\Gamma
\]
是局部常的。等价地，当 $B(G)$ 赋以序拓扑时，映射 $|\Bun_G|\to B(G)$ 是连续的。

\noindent (iv) 对任何 $b\in B(G)$，相应的子函子
\[
  i^b : \Bun_G^b = \Bun_G\times_{|\Bun_G|}\{b\}\subset\Bun_G
\]
是局部闭的，且同构于 $[\ast/\widetilde{G}_b]$，其中 $\widetilde{G}_b$ 是一个群 v-层，使得 $\widetilde{G}_b\to\ast$ 可由局部空间菱形表示，且 $\pi_0\widetilde{G}_b=G_b(E)$。单位元的连通分支 $\widetilde{G}_b^\circ\subset\widetilde{G}_b$ 上同调光滑，维数为 $\langle2\rho,\nu_b\rangle$。

\noindent (v) 特别地，半稳定轨迹 $\Bun_G^{\mathrm{ss}}\subset\Bun_G$ 是开的，且由
\[
  \Bun_G^{\mathrm{ss}} \cong \bigsqcup_{b\in B(G)_{\mathrm{basic}}}[\ast/G_b(E)]
\]
给出。

\noindent (vi) 对任何 $b\in B(G)$，存在映射
\[
  \pi_b : \mathcal{M}_b\to\Bun_G,
\]
它可由局部空间菱形表示、部分紧合且上同调光滑，其中 $\mathcal{M}_b$ 参数化 $G$-丛 $\mathcal{E}$ 连同一个递增的 $\QQ$-滤过，该滤过的分次在所有几何点处同构于 $\mathcal{E}_b$ 及其斜率分次。v-叠 $\mathcal{M}_b$ 可由局部空间菱形表示，且在 $[\ast/G_b(E)]$ 上部分紧合、上同调光滑。

\noindent (vii) v-叠 $\Bun_G$ 是维数为 $0$ 的上同调光滑 Artin 叠。
\end{theorem}

作为例子，我们来分析 $\mathrm{GL}_1$ 与 $\mathrm{GL}_2$ 的情形。对 $\mathrm{GL}_1$ 以及一般的环面，一切都是半稳定的，因此
\[
  \mathrm{Pic} := \Bun_{\mathrm{GL}_1}\cong\bigsqcup_{\ZZ}[\ast/\ul{E}^\times].
\]
对 $\mathrm{GL}_2$，Kottwitz 不变量给出一个分解
\[
  \Bun_{\mathrm{GL}_2} = \bigsqcup_{\alpha\in\frac12\ZZ}\Bun_{\mathrm{GL}_2}^\alpha.
\]
每个连通分支有唯一的半稳定点，由满足 $\kappa(b)=\alpha$ 的基本元 $b\in B(\mathrm{GL}_2)_{\mathrm{basic}}$ 给出。对 $b\in B(\mathrm{GL}_2)_{\mathrm{basic}}\cong\frac12\ZZ$，相应的群 $G_b(E)$ 当 $b\in\ZZ$ 时为 $\mathrm{GL}_2(E)$，而当 $b\in\frac12\ZZ\setminus\ZZ$ 时为 $D^\times$，其中 $D|E$ 是四元数代数。

$\Bun_{\mathrm{GL}_2}$ 的非半稳定点由线丛的扩张给出，它们形如 $O(i)\oplus O(j)$，其中 $i,j\in\ZZ$，且 $2\alpha=i+j$。让我们理解 $\Bun_{\mathrm{GL}_2}$ 内部最简单的退化，即从 $O(\frac12)$ 到 $O\oplus O(1)$ 的退化。这里各层是
\[
  [\ast/D^\times],\quad [\ast/\Aut(O\oplus O(1))].
\]
这里
\[
  \Aut(O\oplus O(1)) = \begin{pmatrix} E^\times & BC(O(1))\\ 0 & E^\times\end{pmatrix}.
\]
\srcpageinl{22}这里 $\mathrm{BC}(\cO(1))$ 可由完满开单位圆盘 $\mathrm{Spd}\,\ol{\FF}_q[[t^{1/p^\infty}]]$ 表示。

在这种情形下，$\Bun_{\mathrm{GL}_2}$ 的局部图卡 $\mathcal{M}_b$ 参数化写成扩张
\[
  0\to\mathcal{L}\to\mathcal{E}\to\mathcal{L}'\to 0
\]
的秩 $2$ 丛 $\mathcal{E}$，使得在所有几何点处都有 $\mathcal{L}\cong\cO$ 与 $\mathcal{L}'\cong\cO(1)$。固定这样的同构定义了一个 $\ul{E}^\times\times\ul{E}^\times$-主丛
\[
  \widetilde{\mathcal{M}}_b\to\mathcal{M}_b,
\]
其中 $\widetilde{\mathcal{M}}_b=\mathrm{BC}(\cO(-1)[1])$ 是一个「负 Banach–Colmez 空间」。这个局部图卡表明，$\Bun_G$ 的局部结构与负 Banach–Colmez 空间的结构密切相关。它也表明，虽然 $\Bun_G$ 的几何相当非标准，它归根到底仍是有限维的、且「光滑」的情形。

对一般的 $G$，我们仍有连通分支的分解
\[
  \Bun_G = \bigsqcup_{\alpha\in\pi_1(G)_\Gamma}\Bun_G^\alpha,
\]
而每个连通分支 $\Bun_G^\alpha$ 都有唯一的半稳定点。

根据 Viehmann 最近的结果 [Vie21]，映射 $|\Bun_G|\to B(G)$ 是同胚。此前这对 $G=\mathrm{GL}_n$ 已由 Hansen [Han17]（基于 [BFH+22]）证明；该论证在 Hamann 未发表的工作中被推广到某些典型群。

关于定理 I.4.1 的证明，我们略作说明。第 (i) 部分本质上已在 [SW20] 中证明，而第 (ii) 部分由 Fargues 与 Anschütz 的结果（定理 I.2.6）推出。在第 (iii) 部分中，关于 $\nu$ 的陈述经 Rapoport–Richartz [RR96] 的论证归约到 $\mathrm{GL}_n$，在那里它就是定理 I.3.4。关于 $\kappa$ 的陈述则需要更多工作，至少一般情形下如此：若 $G$ 的导群单连通，则可归约到环面，而环面的情形不难处理。一般地，一种办法是用 z-扩张 $\widetilde{G}\to G$ 归约到导群单连通的情形。为此需要 $\Bun_{\widetilde{G}}\to\Bun_G$ 是 v-叠的满射；我们用 Beauville–Laszlo 单值化来证明这一点。另一种办法是使用 Borovoi [Bor98] 的阿贝尔化 Kottwitz 集，我们证明它在完满空间 $S$ 上相对地表现良好。第 (iv) 部分也是定理 I.3.4 的推论\yizhu{原文此处作「is a also consequence」，多出一个冠词 a，属排版笔误，译文按「也是……的推论」处理。}{A：原文 p.22}
。第 (v) 部分是第 (iii) 与 (iv) 部分的推论。于是关键落在第 (vi) 部分，它将形式地蕴含 (vii)。（人们也可以、并且我们确实也直接证明第 (vii) 部分。然而第 (vi) 部分对下文所证的其他结果至关重要。）$\mathcal{M}_b$ 本身的性质容易建立——上面关于 $\mathrm{GL}_2$ 的分析容易推广，以说明 $\widetilde{\mathcal{M}}_b$ 是负 Banach–Colmez 空间的逐次扩张。关键的困难在于证明
\[
  \pi_b : \mathcal{M}_b\to\Bun_G
\]
是上同调光滑的。请注意，由于我们是在完满空间中工作，这里没有切空间，因而不能指望通过形变理论来证明光滑性。为处理这个问题，我们仍然证明了一般性的「上同调光滑性的 Jacobian 判别法」。这里的设定如下。

设 $S$ 是完满空间，$Z\to X_S$ 是（sousperfectoid）adic 空间的光滑映射；这意味着 $Z$ 是一个在 $X_S$ 上某个有限维球之上局部 \text{ét}ale 的 adic 空间。在这种情形下，我们可以定义一个 v-层 $\mathcal{M}_Z\to S$，参数化 $Z\to X_S$ 的截面，即：对 $S'/S$ 为完满空间，其 $S'$-值点由提升 $X_{S'}\to X_S$ 的映射 $s: X_{S'}\to Z$ 给出。对于
